% TOP_PROBLEM: {"id": 30003584, "title": "Nuclear Dimension of Transformation-Group Crossed Products", "category_id": 4, "category": {"id": 4, "name": "algebra", "display_name": "Algebra", "description": "Group theory, ring theory, field theory, and algebraic structures.", "slug": "algebra", "order_index": 4, "created_at": "2026-07-31T15:26:25.671Z"}, "status": "partially_solved"}
% FIRST_POSTED: 2026-10-06
% ATTEMPT_STATUS: unresolved
% OPENAI_CLAIM: {"schema_version":1,"source_repo":"https://github.com/openai/math","source_commit":"adc7f1241b42e322a6451854ab7e4b4c146bf78a","release_date":"2026-10-06","match":"partial","resolution":"partial","independently_reviewed":false,"second_pass":"n/a","families":[{"family":"302","title":"Radius of comparison equals half the mean dimension","subject":"Operator algebras","manuscripts":[{"title":"Radius of comparison equals half the mean dimension","pdf_path":"preprints/Radius-of-comparison-equals-half-the-mean-dimension-September-25-2026/paper.pdf","readme_path":"preprints/Radius-of-comparison-equals-half-the-mean-dimension-September-25-2026/README.md","date":"2026-09-25","theorem_ref":"thm:main (sections/01-introduction.tex:51)","statement_sources":[{"path":"preprints/Radius-of-comparison-equals-half-the-mean-dimension-September-25-2026/build/sections/01-introduction.tex","line":51,"label":"thm:main"}]}],"lean":null,"reasoning_trace":null}]}
\documentclass[11pt]{article}
\usepackage[margin=25mm]{geometry}
\usepackage{fontspec}
\setmainfont{Latin Modern Roman}
\usepackage{amsmath,amssymb,mathtools}
\usepackage{unicode-math}
\setmathfont{Latin Modern Math}
\usepackage{enumitem}
\usepackage{xurl,hyperref,bookmark}
\hypersetup{colorlinks=true,urlcolor=blue}
\setcounter{secnumdepth}{0}
\setlength{\parindent}{0pt}
\setlength{\parskip}{5pt}
\setlength{\emergencystretch}{3em}
\begin{document}
\section*{Nuclear Dimension of Transformation-Group Crossed Products}

\subsection{OpenAI's claim}
Whenever an infinite compact metrizable system is minimal under a homeomorphism, its crossed-product comparison radius is claimed to be half its mean dimension, including infinite values; zero forces nuclear dimension at most one. The cited manuscript argues for that result; its relation to this problem is specified below.

\subsection{Mathematical statement of the claimed result}
\paragraph{Mathematical claim 1.} Let $h$ be a homeomorphism of $X$ whose orbits are all dense, where $X$ is compact, infinite and metrizable. OpenAI claims
\[
\operatorname{rc}(C(X)\rtimes_h\mathbb Z)=\tfrac12\operatorname{mdim}(X,h)\quad\text{in }[0,\infty].
\]
For positive stabilized elements define $d_\tau(a)=\lim_m\tau(a^{1/m})$. The radius of comparison is the infimum of $r\ge0$ such that $d_\tau(a)+r<d_\tau(b)$ for every normalized trace implies Cuntz subequivalence $a\precsim b$. Mean dimension is the asymptotic covering dimension per iterate, defined from joins of finite open covers. Both finite and infinite mean dimensions are included.

\subsection{Outline of the argument}
Manuscript 1 (family 302, 2026-09-25): Finite-support comparison witnesses are converted to bundle embeddings, which feed the companion's boundary-preserving compression theorem. Packet reconstruction translates that compression into a smaller cover without introducing absent colors. Applying the construction on aligned orbit blocks contradicts any larger cover-growth rate, providing the missing lower bound; established upper bounds and regularity results finish the equality and zero-dimensional consequences.

\subsection{Scope relative to this problem}
Manuscript 1 (family 302, 2026-09-25): For minimal Z actions on infinite compact metrizable spaces, rc=mdim/2 and the zero-dimension corollary determine finite nuclear dimension; this problem asks a general locally compact-group transformation-group classification, including nonminimal and nonmetrizable cases.

\subsection{Formal verification}
Family 302 has no Lean scope page in this snapshot; no formalization of these manuscript results is listed.

This site has not executed the Lean files or checked their dependencies. These descriptions report the scope of the linked declarations, rather than a new verification of the proofs.

\subsection{Status caveat}
OpenAI attributes these results to an internal model. Its release notice cautions that results lacking full formalization can contain mistakes. This site has not independently assessed the proofs; a solved label records the scope of the reported claim.

\subsection{Sources}
\begin{itemize}
\item OpenAI, \emph{Radius of comparison equals half the mean dimension}, 2026-09-25; citation key \texttt{OAI:Radius-of-comparison-equals-half-the-mean-dimension-September-25-2026}. \href{https://github.com/openai/math/blob/adc7f1241b42e322a6451854ab7e4b4c146bf78a/preprints/Radius-of-comparison-equals-half-the-mean-dimension-September-25-2026/paper.pdf}{PDF (GitHub)}; \href{https://github.com/openai/math/raw/adc7f1241b42e322a6451854ab7e4b4c146bf78a/preprints/Radius-of-comparison-equals-half-the-mean-dimension-September-25-2026/paper.pdf}{PDF (direct)}; \href{https://github.com/openai/math/blob/adc7f1241b42e322a6451854ab7e4b4c146bf78a/preprints/Radius-of-comparison-equals-half-the-mean-dimension-September-25-2026/README.md}{Manuscript page and citation}; \href{https://github.com/openai/math/blob/main/preprints/Radius-of-comparison-equals-half-the-mean-dimension-September-25-2026/paper.pdf}{Latest version}; Statement source: \href{https://github.com/openai/math/blob/adc7f1241b42e322a6451854ab7e4b4c146bf78a/preprints/Radius-of-comparison-equals-half-the-mean-dimension-September-25-2026/build/sections/01-introduction.tex\#L51}{thm:main}.
\item \href{https://github.com/openai/math/blob/adc7f1241b42e322a6451854ab7e4b4c146bf78a/CONTENTS.md}{OpenAI catalogue entry}.
\end{itemize}
\end{document}
